\section{Counting zeros with the Thorin measure}\label{ch:turing}
\label{sec:turing}


Turing's method certifies the count $N(T)$ without evaluating $\zeta$ off the critical line, from the sign changes of $Z$ and an explicit bound on the integral of $S$, where $\pi S(t)=\arg\zeta(\tfrac12+it)$; Backlund's method, which it replaced, bounds the argument of $\zeta$ along a horizontal segment off the line (Section~\ref{sec:ch5:turing}). Without the assumption of Proposition~\ref{prop:mass} that the zeros below the window lie on the line, the argument principle gives at every resolution an exact identity for $N(T)$ and for the number $\nu(\tfrac12+\varepsilon,T)$ of zeros to the right of the line of evaluation, so that off line zeros are counted and not only detected. At $\varepsilon=\tfrac12$ the identity is Backlund's count with the vertical side on $\re s=1$.

\subsection{The cumulative Thorin mass as a smoothed counting function}\label{sec:ch9:mass}

Recall that $M_\varepsilon(\theta)=\mu_\varepsilon((0,\theta^2])=\pi^{-1}\arg\xi(\tfrac12+\varepsilon+i\theta)$, with the argument continued along the line $\re s=\tfrac12+\varepsilon$ from $\theta=0$, where $\xi>0$ (Sections~\ref{sec:ch6:consequences} and \ref{ch:distances}). Since $\xi(s)=\tfrac12s(s-1)\pi^{-s/2}\Gamma(s/2)\zeta(s)$, the mass splits as
\[
M_\varepsilon(\theta)=N^{\rm sm}_\varepsilon(\theta)+S_\varepsilon(\theta),\qquad S_\varepsilon(\theta)=\frac1\pi\arg\zeta\Big(\frac12+\varepsilon+i\theta\Big),
\]
where, with $s=\tfrac12+\varepsilon+i\theta$ and $\theta>0$, the smooth part $N^{\rm sm}_\varepsilon(\theta)=\pi^{-1}\big[\arg s+\arg(s-1)+\im\log\Gamma(\tfrac14+\tfrac\varepsilon2+\tfrac{i\theta}2)-\tfrac\theta2\log\pi\big]$ comes from $s(s-1)\pi^{-s/2}\Gamma(s/2)$, with the principal arguments of $s$ and $s-1$, which lie in $(0,\pi)$, and the branch of $\log\Gamma$ that is real on the positive axis. Thus $S_\varepsilon$ is $\pi^{-1}\arg\zeta$ continued along the vertical line from $S_\varepsilon(0+)=-1$ if $0<\varepsilon<\tfrac12$, when $\arg(s-1)\to\pi$ as $\theta\downarrow0$, and from $S_\varepsilon(0+)=0$ if $\varepsilon>\tfrac12$.

As $\varepsilon\downarrow0$, $N^{\rm sm}_\varepsilon\to\vartheta/\pi+1$, where $\vartheta(\theta)=\im\log\Gamma(\tfrac14+\tfrac{i\theta}2)-\tfrac\theta2\log\pi$ is the Riemann Siegel theta function, and for $\theta$ not an ordinate Proposition~\ref{prop:zerodensity} below gives $M_\varepsilon(\theta)\to N(\theta)-2\nu(\tfrac12,\theta)$, where $\nu(\sigma,\theta)$ counts the zeros with real part greater than $\sigma$ and ordinate in $(0,\theta]$ (Section~\ref{ch:intro}). When no zero below height $\theta$ lies to the right of the critical line, the limit is therefore the Riemann von Mangoldt formula $N(\theta)=\vartheta(\theta)/\pi+1+S(\theta)$ \cite[Section 9.3]{Titchmarsh}.

Integrating Proposition~\ref{prop:struct} in $\theta$ writes $M_\varepsilon$ as a sum of Cauchy steps. Under RH, for $\theta\ge0$,
\begin{gather*}
M_\varepsilon(\theta)=\sum_{\gamma>0}\big[F_\varepsilon(\theta-\gamma)-F_\varepsilon(-\gamma)+F_\varepsilon(\theta+\gamma)-F_\varepsilon(\gamma)\big],\\
F_\varepsilon(x)=\frac12+\frac1\pi\arctan\frac x\varepsilon ,
\end{gather*}
with the sum over all positive ordinates. Since $F_\varepsilon(-\gamma)+F_\varepsilon(\gamma)=1$, this says that $M_\varepsilon=P_\varepsilon*\tilde N$, where $\tilde N$ is the odd extension of $N$. As $\vartheta$ is odd, $S_\varepsilon$ is then the Poisson average of the odd extension of $S$ up to an explicit smooth function. If only the zeros below height $T+B$ are known to lie on the line, the display holds on $[0,T]$ up to an error of order $\varepsilon\delta^2/B^3$ for each off line pair above $T+B$ at distance $\delta$ from the line, counted as a double ordinate, because $P_{\varepsilon+\delta}(x)+P_{\varepsilon-\delta}(x)-2P_\varepsilon(x)=-6\varepsilon\delta^2/(\pi x^4)+O(x^{-6})$.

For a single step,
\begin{align*}
\int_{-L}^{L}\big|F_\varepsilon(x)-\mathbf 1_{x\ge0}\big|\,dx&=\frac2\pi\Big[L\arctan\frac\varepsilon L+\frac\varepsilon2\log\Big(1+\frac{L^2}{\varepsilon^2}\Big)\Big]\\
&=\frac{2\varepsilon}\pi\Big(\log\frac L\varepsilon+1\Big)+O\Big(\frac{\varepsilon^3}{L^2}\Big),
\end{align*}
which explains the $\varepsilon\log(1/\varepsilon)$ rate of the $L^1$ distance in Table~\ref{tab:tv}. Taking half of it on each side of each of the $79$ zeros below height $200$, with $L$ the half gap on that side, gives $1.61$ at $\varepsilon=0.005$, against the computed $1.55$; the difference comes from the tails of the other steps.

\subsection{A Littlewood identity and the zero density}\label{sec:ch9:density}

\begin{proposition}\label{prop:littlewood}
Let $T>0$. For every $\varepsilon>0$ such that $\xi$ has no zero on the segment from $\tfrac12+\varepsilon$ to $\tfrac12+\varepsilon+iT$,
\[
\pi M_\varepsilon(T)=\frac{\partial}{\partial\varepsilon}\int_0^T\log\Big|\xi\Big(\frac12+\varepsilon+it\Big)\Big|\,dt,
\]
and for every $\varepsilon>0$
\[
\int_0^T\log\Big|\xi\Big(\frac12+\varepsilon+it\Big)\Big|\,dt=\int_0^T\log\Big|\xi\Big(\frac12+it\Big)\Big|\,dt+\pi\int_0^\varepsilon M_{\varepsilon'}(T)\,d\varepsilon' .
\]
\end{proposition}

\begin{proof}
Since $\partial_\sigma\log|\xi(\sigma+it)|=\re(\xi'/\xi)(\sigma+it)=v_\sigma(t)$ and $v_\sigma(t)=\partial_t\arg\xi(\sigma+it)$, we have $M_\varepsilon(T)=\pi^{-1}\int_0^Tv_{1/2+\varepsilon}(t)\,dt$, and when no zero lies on the segment, differentiation under the integral gives the first identity. The integral $I(\varepsilon)$ on the left of the second is continuous on $[0,\infty)$, because the logarithmic singularities at the zeros are integrable uniformly in $\varepsilon$, and continuously differentiable except at the finitely many $\varepsilon'$ at which the segment meets a zero. There $I'$ jumps by $2\pi$ for each zero met at an interior point of the segment, since $\int_{-c}^{c}a\,du/(a^2+u^2)\to\pm\pi$ as $a\to0\pm$. So $I$ is absolutely continuous, and $I(\varepsilon)-I(0)=\int_0^\varepsilon I'(\varepsilon')\,d\varepsilon'$ is the second identity.
\end{proof}

The jump of $2\pi$ in the slope $\pi M_\varepsilon(T)$ when the line crosses a zero is the term $2\nu$ of Proposition~\ref{prop:zerodensity} below. Littlewood's lemma \cite[Section 9.9]{Titchmarsh} is the same Cauchy Riemann identity with the argument taken along a path through $\re s=2$; the two arguments differ by $2\pi\nu(\tfrac12+\varepsilon,T)$, and Exercise~\ref{ex:ch9:littlewood} derives the lemma from the two propositions. Let $\tau_\varepsilon(T)=\pi^{-1}\Delta\arg\xi$ along the horizontal segment from $\tfrac12+\varepsilon+iT$ to $\tfrac12+iT$.

\begin{proposition}\label{prop:zerodensity}
Let $\varepsilon>0$ and $T>0$, with $T$ not an ordinate and no zero of $\xi$ on the segment from $\tfrac12+\varepsilon$ to $\tfrac12+\varepsilon+iT$. Then
\[
M_\varepsilon(T)+\tau_\varepsilon(T)=N(T)-2\,\nu\Big(\frac12+\varepsilon,T\Big).
\]
The segment condition holds at $\varepsilon=\tfrac12$, and $\nu(1,T)=0$, which gives the unconditional count $N(T)=M_{1/2}(T)+\tau_{1/2}(T)$, so that $\nu(\tfrac12+\varepsilon,T)=\tfrac12\big[N(T)-M_\varepsilon(T)-\tau_\varepsilon(T)\big]$ for every $\varepsilon$ satisfying the hypotheses.
\end{proposition}

\begin{proof}
Let $\sigma_1>1$. Write $\Delta_{\rm right}$ for the change of $\arg\xi$ along the segment from $\sigma_1$ to $\sigma_1+iT$, and $\Delta_{\rm top}(a\to b)$ for its change along the segment from $a+iT$ to $b+iT$. The argument principle on the rectangle $[\tfrac12+\varepsilon,\sigma_1]\times[0,T]$, on whose bottom edge $\xi$ is positive, gives $\Delta_{\rm right}+\Delta_{\rm top}(\sigma_1\to\tfrac12+\varepsilon)-\pi M_\varepsilon(T)=2\pi\nu(\tfrac12+\varepsilon,T)$. Subtracting this from the classical formula $\pi N(T)=\Delta_{\rm right}+\Delta_{\rm top}(\sigma_1\to\tfrac12)$ \cite[Section 9.3]{Titchmarsh} leaves $\Delta_{\rm top}(\tfrac12+\varepsilon\to\tfrac12)+\pi M_\varepsilon(T)=\pi N(T)-2\pi\nu(\tfrac12+\varepsilon,T)$, and $\Delta_{\rm top}(\tfrac12+\varepsilon\to\tfrac12)=\pi\tau_\varepsilon(T)$. At $\varepsilon=\tfrac12$ the segment lies on $\re s=1$, where $\zeta$ has no zeros, and $\nu(1,T)=0$ because the Euler product has no zeros in $\re s>1$.
\end{proof}

Half the difference of $M_\varepsilon(T)+\tau_\varepsilon(T)$ at two resolutions counts the zeros in the strip between the two lines. The factor two comes from the functional equation: a zero $\tfrac12+\delta+i\gamma$ with $\delta>\varepsilon$ and its partner $\tfrac12-\delta+i\gamma$ add $2$ to $N(T)$ and $1$ to $\nu(\tfrac12+\varepsilon,T)$, hence nothing to $M_\varepsilon(T)+\tau_\varepsilon(T)$, whereas two zeros on the line at the same ordinate add $2$. In the density the pair enters through the kernels $P_{\varepsilon+\delta}$ and $-P_{\delta-\varepsilon}$ of Proposition~\ref{prop:struct}, whose masses cancel, as in the proof of Proposition~\ref{prop:signature}(iii) (Exercise~\ref{ex:ch9:pair}).

Zero density estimates such as Ingham's \cite{Ingham40} bound the number $N(\sigma,T)$ of zeros with real part at least $\sigma$ and ordinate in $(0,T]$, which is at least $\nu(\sigma,T)$, by $T^{e(\sigma)+o(1)}$, with an exponent $e(\sigma)<1$ for every $\sigma>\tfrac12$, and Section~\ref{ch:almost} turns these bounds, with no hypothesis, into a bound on the mass of the Cauchy kernels of the zeros to the right of the basepoint $\sigma$, the negative part of the split \eqref{ag:eq:split} (Theorem~\ref{ag:thm:almost}). Proposition~\ref{prop:zerodensity} evaluates $\nu$ itself, as an integer, but only up to a finite height.

\subsection{The count and the zero density below height 200}\label{sec:ch9:count}

The right side of the identity is an integer, so a computation of the left side with error below one half determines it, provided no winding of the argument is missed. We follow the argument on a grid, summing the principal values of the arguments of quotients of successive values of $\xi$, with steps $\varepsilon/8$ ($\varepsilon/4$ at $\varepsilon=0.005$) up the vertical side from $\theta=0$ and $4000$ steps along the horizontal segment. A certificate would need, besides ball arithmetic at the grid points, a bound on the change of the argument between them. We do not supply that bound, and the conclusions below are computed, not certified.

Table~\ref{tab:zerodensity} gives the identity at $T=200$, evaluated by the P\'olya route. Every row returns $79$ to within $6\times10^{-13}$. The row $\varepsilon=\tfrac12$ gives $N(200)=79$, and with the $79$ sign changes of $\Xi$ on $(0,200)$ (Section~\ref{sec:ch4:height200}) this places every zero below height $200$ on the line. Independently of the sign changes, the rows $\varepsilon<\tfrac12$ give $\nu(\tfrac12+\varepsilon,200)=0$ for every $\varepsilon\ge0.005$, since $\nu$ is nonincreasing in $\sigma$, so that by the functional equation no zero below height $200$ lies farther than $0.005$ from the line.

\begin{table}[htbp]
\centering\small
\begin{tabular}{rrrr}
\toprule
$\varepsilon$ & $M_\varepsilon(200)$ & $\tau_\varepsilon(200)$ & $M_\varepsilon+\tau_\varepsilon=N-2\nu$\\
\midrule
1 & 79.36737810665 & $-0.36737810665$ & 79.00000000000003\\
0.5 & 79.19416875933 & $-0.19416875933$ & 78.99999999999993\\
0.25 & 79.09886991193 & $-0.09886991193$ & 79.00000000000024\\
0.1 & 79.03976768776 & $-0.03976768776$ & 79.00000000000006\\
0.05 & 79.01989993142 & $-0.01989993142$ & 78.99999999999967\\
0.02 & 79.00796178162 & $-0.00796178162$ & 78.99999999999970\\
0.01 & 79.00398102016 & $-0.00398102016$ & 78.99999999999993\\
0.005 & 79.00199052625 & $-0.00199052625$ & 78.99999999999947\\
\bottomrule
\end{tabular}
\caption{Proposition~\ref{prop:zerodensity} at $T=200$ by the P\'olya route (computed, not certified).}
\label{tab:zerodensity}
\end{table}

The identity also gives the mass defect of Proposition~\ref{prop:mass} exactly. If $T$ is not an ordinate and every zero with ordinate in $(0,T]$ lies on the line, then $\nu(\tfrac12+\varepsilon,T)=0$ and $M_\varepsilon(T)=N(T)-\tau_\varepsilon(T)$ for every $\varepsilon>0$. Comparing with that proposition and Remark~\ref{rem:ch7:cubic} gives $\tau_\varepsilon(T)=(\varepsilon/\pi)\,\Xi'(T)/\Xi(T)+O(\varepsilon^3)$. At $T=200$ the first order term is $-0.3981063\,\varepsilon$, minus the mass defect found in Section~\ref{sec:ch7:cumulative}, and for $\varepsilon\le0.1$ the column $\tau_\varepsilon(200)$ of Table~\ref{tab:zerodensity} exceeds it by about $0.043\,\varepsilon^3$.
